% Distribution metrics, FoodLogAthl-218 vs UEC-Food 256. Values are the supplied ones (table_xj.tex); layout transposed
% to one column (metrics as rows) and the broken \cmidrule{11-13} of the input removed.
% Arrows give the direction that indicates a more diverse / harder-to-separate dataset (see caption).
\begin{table}[tb]
  \centering
  \caption{Feature-space distribution metrics for UEC-Food 256 and FoodLogAthl-218 under three feature extractors.
  The arrow gives the direction that indicates greater spread and less separable classes; the more diverse value of each
  pair is in bold. Trace and Davies--Bouldin: higher means greater spread; Silhouette and Calinski--Harabasz: lower means
  less compact, more overlapping classes. Single run per cell, no confidence intervals.}
  \label{table:diversity_metrics_comparison}
  \small
  \setlength{\tabcolsep}{3pt}
  \resizebox{\columnwidth}{!}{\begin{tabular}{@{}l ccc ccc@{}}
    \toprule
    & \multicolumn{3}{c}{UEC-Food 256} & \multicolumn{3}{c}{FoodLogAthl-218 (ours)} \\
    \cmidrule(lr){2-4} \cmidrule(lr){5-7}
    Metric & CLIP & EVA-CLIP & Swin & CLIP & EVA-CLIP & Swin \\
    \midrule
    Trace ($\uparrow$)               & 71.21  & 45.09  & 264.00  & \textbf{81.70}  & \textbf{49.92}  & \textbf{282.57} \\
    Silhouette ($\downarrow$)        & 0.0335 & 0.0577 & $-$0.0064 & \textbf{$-$0.0210} & \textbf{0.0036} & \textbf{$-$0.0288} \\
    Davies--Bouldin ($\uparrow$)     & 3.64   & 3.02   & 4.11    & \textbf{4.52}   & \textbf{3.68}   & \textbf{4.51}  \\
    Calinski--Harabasz ($\downarrow$)& 74.05  & 115.80 & 49.17   & \textbf{27.51}  & \textbf{51.51}  & \textbf{21.29} \\
    \bottomrule
  \end{tabular}}
\end{table}
